\section{Mathematical framework}
\label{sec:math}

The direct conversion consists of three separate mathematical layers:
\begin{enumerate}
  \item compress a halo mass history into \diffmah{} parameters;
  \item compress a stellar CoG into five readable, ordered coordinates;
  \item learn a conditional probability distribution between those vectors.
\end{enumerate}
Keeping these layers separate is essential. A low profile-reconstruction error
does not imply that halo properties can predict the coordinates, and a good
conditional mean does not imply that the model reproduces the galaxy population.

\subsection{The adopted \diffmah{} parameterization}
\label{sec:diffmah}

\diffmah{} describes halo growth as a rolling power law in cosmic time
\citep{Hearin2021}. Let $u=\logten(t/\gyr)$ and define the logistic function
$S(a)=(1+e^{-a})^{-1}$. The time-dependent logarithmic growth index is
\begin{equation}
  \alpha(u)=\alpha_{\rm early}+
  \left(\alpha_{\rm late}-\alpha_{\rm early}\right)
  S\!\left[k(u-u_c)\right],
  \qquad k=3.5,
  \label{eq:diffmah_alpha}
\end{equation}
and the implemented mass history is
\begin{equation}
  \logten \Mpeak(t)=\logten M_p+\alpha(u)(u-u_0).
  \label{eq:diffmah_mass}
\end{equation}
Here $u_0=\logten(t_0/\gyr)$ is the fitting endpoint. The model passes through
$\logten M_p$ at $t=t_0$. The four fitted halo coordinates are
\begin{equation}
  \vect{\theta}_{\rm MAH}=
  \left(\logten M_p,\,u_c,\,\alpha_{\rm early},\,\alpha_{\rm late}\right).
  \label{eq:mah_vector}
\end{equation}

\begin{table}[htbp]
  \centering
  \caption{Halo-side parameters used by the direct conversion. ``Early'' and
  ``late'' refer to the asymptotic logarithmic growth indices, not stellar
  formation rates.}
  \label{tab:halo_parameters}
  \begin{tabularx}{\textwidth}{@{}l l X@{}}
    \toprule
    Symbol & Code name & Meaning \\
    \midrule
    $\logten M_p$ & \texttt{logmp} & Peak-mass normalization in $\logten(\Msun)$ at the chosen endpoint $t_0$ \\
    $u_c=\logten(t_c/\gyr)$ & \texttt{logtc} & Cosmic time around which the growth index transitions \\
    $\alpha_{\rm early}$ & \texttt{early} & Fast-regime logarithmic growth index \\
    $\alpha_{\rm late}$ & \texttt{late} & Slow-regime logarithmic growth index \\
    $\ctwo$ & \texttt{c200c} & Halo concentration at a stated epoch \\
    \bottomrule
  \end{tabularx}
\end{table}

The four numbers depend on the time range and endpoint used in the fit. In the
descendant-conditioned model, $t_0=t(z=0.4)$ and the complete history through
that time is fitted. In the epoch-local model, $t_0=t(z)$ and only history up to
the prediction epoch is fitted. Thus an epoch-local $\logten M_p$ is close to
the concurrent peak mass, while the descendant-conditioned normalization is a
final-descendant quantity.

\subsection{Readable CoG coordinates in \ExpFifty{}}
\label{sec:exp50_coordinates}

Define the central fraction
\begin{equation}
  f_2\equiv F_\star(2\kpc)
  =\frac{M_\star(<2\kpc)}{M_{\star,\rm tot}},
  \label{eq:f2}
\end{equation}
and ordinary enclosed-mass radii by
\begin{equation}
  F_\star(R_p)=p,\qquad p\in\{0.2,0.5,0.8\}.
  \label{eq:ordinary_quantiles}
\end{equation}
To guarantee $R_{20}<R_{50}<R_{80}$ for every mean prediction and stochastic
draw, \ExpFifty{} models positive gaps in logarithmic radius. With
$r_p=\logten(R_p/\kpc)$, its coordinate vector is
\begin{equation}
  \vect{\theta}_{\cog}^{(50)}=
  \begin{pmatrix}
    \logten(M_{\star,\rm tot}/\Msun)\\
    \operatorname{logit}(f_2)\\
    r_{20}\\
    \logten(r_{50}-r_{20})\\
    \logten(r_{80}-r_{50})
  \end{pmatrix}.
  \label{eq:exp50_coordinates}
\end{equation}
The final two coordinates are logarithms of separations in $\logten R$; they
are not logarithms of $R_{50}-R_{20}$ in linear kiloparsecs. Decoding exponentiates
the gaps and cumulatively adds them to $r_{20}$.

\subsection{Central-censored CoG coordinates in \ExpFiftyOne{}}
\label{sec:exp51_coordinates}

At high redshift, a valid compact galaxy can already have $f_2>0.5$. Its
ordinary $R_{20}$ and $R_{50}$ then lie below the innermost measured aperture
and cannot be inferred from the CoG grid. \ExpFiftyOne{} retains $f_2$ as an
explicit coordinate and defines quantiles only for the mass outside 2 kpc. For
$q\in\{0.2,0.5,0.8\}$, define the corresponding fraction of total mass
\begin{equation}
  p_q^{\rm out}=f_2+(1-f_2)q,
  \label{eq:outer_level}
\end{equation}
and the outer-mass radius by
\begin{equation}
  F_\star(R_{q,\rm out})=p_q^{\rm out}.
  \label{eq:outer_quantiles}
\end{equation}
The coordinate vector is then
\begin{equation}
  \vect{\theta}_{\cog}^{(51)}=
  \begin{pmatrix}
    \logten(M_{\star,\rm tot}/\Msun)\\
    \operatorname{logit}(f_2)\\
    r_{20,\rm out}\\
    \logten(r_{50,\rm out}-r_{20,\rm out})\\
    \logten(r_{80,\rm out}-r_{50,\rm out})
  \end{pmatrix}.
  \label{eq:exp51_coordinates}
\end{equation}
Unlike assigning a boundary value to an unresolved ordinary radius, this
definition is measured for every retained profile and preserves the compact
galaxy as scientific information.

\begin{figure}[htbp]
  \centering
  \resizebox{0.98\textwidth}{!}{%
  \begin{tikzpicture}[
      font=\small,
      node distance=8mm and 11mm,
      box/.style={draw,rounded corners,align=center,minimum height=13mm,minimum width=34mm,fill=white},
      arrow/.style={-{Latex[length=2.3mm]},thick},
      note/.style={align=center,font=\footnotesize,text=hsgray}
    ]
    \node[box,draw=hsblue,very thick] (mah) {$\Mpeak(t)$\\main-branch history};
    \node[box,draw=hsblue,very thick,right=of mah] (dmah) {$\vect{\theta}_{\rm MAH}$\\four \diffmah{} parameters};
    \node[box,draw=hsgreen,very thick,right=of dmah] (map) {$P(\vect{\theta}_{\cog}\mid\vect{x})$\\mean + covariance};
    \node[box,draw=hsorange,very thick,right=of map] (coords) {$\vect{\theta}_{\cog}$\\five readable coordinates};
    \node[box,draw=hsorange,very thick,right=of coords] (cog) {$M_\star(<R)$\\monotone CoG};
    \draw[arrow] (mah) -- (dmah);
    \draw[arrow] (dmah) -- (map);
    \draw[arrow] (map) -- (coords);
    \draw[arrow] (coords) -- (cog);
    \node[note,below=4mm of dmah] {history compression};
    \node[note,below=4mm of map] {probabilistic regression};
    \node[note,below=4mm of coords] {profile compression};
    \node[box,draw=hspurple,dashed,below=17mm of map] (conc) {$\ctwo(z)$\\secondary halo state};
    \draw[arrow,dashed] (conc) -- (map);
  \end{tikzpicture}
  }
  \caption{The direct conversion has three independently testable layers. A
  failure in profile compression, halo-to-coordinate prediction, or stochastic
  calibration should not be attributed to either of the other layers.}
  \label{fig:pipeline}
\end{figure}

\subsection{Monotone profile decoder}

For a predicted coordinate vector, the decoder constructs knots
$(\logten R_i,F_i)$ from the 2-kpc anchor, the three quantile radii, and the
outer boundary $(148.2\kpc,1)$. It uses a shape-preserving piecewise cubic
Hermite interpolator \citep{Fritsch1984}, clips the result to $0<F\leq1$, and
enforces cumulative monotonicity. The reconstructed profile is
\begin{equation}
  \logten\!\left[\frac{M_\star(<R)}{\Msun}\right]
  =\logten\!\left(\frac{M_{\star,\rm tot}}{\Msun}\right)
  +\logten F_\star(R).
  \label{eq:decoder}
\end{equation}
This decoder does not assume a Sersic, Moffat, or other global analytic profile
family. It assumes only the information in the five coordinates plus smooth,
monotone interpolation between their knots.

\subsection{Conditional mean and residual distribution}
\label{sec:emulator_math}

Let the halo feature vector be
\begin{equation}
  \vect{x}=(\logten M_p,u_c,\alpha_{\rm early},\alpha_{\rm late},\ctwo),
\end{equation}
and standardize each feature within the training sample,
\begin{equation}
  x_{s,j}=\frac{x_j-\overline{x}_j}{s_{x,j}}.
  \label{eq:standardization}
\end{equation}
The linear conditional mean uses
\begin{equation}
  \vect{\phi}_{\rm lin}(\vect{x}_s)=(1,x_{s,1},\ldots,x_{s,5}).
\end{equation}
The degree-2 model appends the seven terms selected in the earlier emulator
experiments,
\begin{align}
  \vect{\phi}_{2} = \big(&\vect{\phi}_{\rm lin},
  x_{\log M_p}^2,
  x_{u_c}x_{\alpha_{\rm late}},
  x_{\alpha_{\rm early}}^2,
  x_{\alpha_{\rm late}}x_{c},
  x_{\alpha_{\rm late}}^2,\nonumber\\
  &x_{\log M_p}x_{\alpha_{\rm early}},
  x_{\log M_p}x_{\alpha_{\rm late}}\big).
  \label{eq:poly2_basis}
\end{align}
For target coordinate $a$, the mean is
\begin{equation}
  \mu_a(\vect{x})=\vect{\beta}_a^{\mathsf T}\vect{\phi}(\vect{x}_s).
  \label{eq:mean_model}
\end{equation}

Scatter is allowed to vary from halo to halo. Its marginal variance is
\begin{equation}
  \sigma_a^2(\vect{x})=
  \exp\!\left[\vect{\gamma}_a^{\mathsf T}(1,\vect{x}_s)\right],
  \label{eq:heteroscedastic}
\end{equation}
with an $L_2$ penalty on the non-intercept variance slopes. A fixed residual
correlation matrix $\vect{C}$ couples the five coordinates, giving
\begin{equation}
  \vect{\Sigma}(\vect{x})=
  \operatorname{diag}[\vect{\sigma}(\vect{x})]\,
  \vect{C}\,
  \operatorname{diag}[\vect{\sigma}(\vect{x})].
  \label{eq:covariance}
\end{equation}
The working predictive model is therefore
\begin{equation}
  \vect{\theta}_{\cog}\mid\vect{x}
  \sim \mathcal{N}\!\left[\vect{\mu}(\vect{x}),\vect{\Sigma}(\vect{x})\right],
  \label{eq:conditional_gaussian}
\end{equation}
followed by the nonlinear monotone decoder in \cref{eq:decoder}. Gaussianity
applies in transformed coordinate space, not directly to stellar mass at each
radius.

\subsection{Prediction and evaluation statistics}
\label{sec:metrics}

All predictive comparisons are out of fold. For galaxy $i$ with $K=24$ CoG
radii, the conditional-mean profile error is summarized by
\begin{equation}
  \mathrm{RMS}_i=\left[\frac{1}{K}\sum_{k=1}^{K}
  \left(\widehat{m}_{ik}-m_{ik}\right)^2\right]^{1/2},
  \label{eq:profile_rms}
\end{equation}
where $m_{ik}=\logten[M_{\star,i}(<R_k)/\Msun]$. The report quotes the median
of $\mathrm{RMS}_i$ unless otherwise stated.

Probabilistic sharpness and calibration are measured with empirical CRPS
\citep{Gneiting2007}. For truth $y$ and $D$ draws $Y_d$,
\begin{equation}
  \widehat{\mathrm{CRPS}}(y)=
  \frac{1}{D}\sum_{d=1}^{D}|Y_d-y|
  -\frac{1}{2D^2}\sum_{d=1}^{D}\sum_{d'=1}^{D}|Y_d-Y_{d'}|.
  \label{eq:crps}
\end{equation}
The profile CRPS is averaged over galaxies and radii. Lower CRPS is better; it
rewards a distribution that is both concentrated and statistically calibrated.

Population comparisons use energy distance \citep{Szekely2013}. Because a
finite TNG sample compared with a resampled copy of itself has nonzero energy
distance, results are divided by this finite-sample floor. A ratio near unity
means the model population is statistically as close to TNG as two finite
realizations of the TNG population are to each other.
